\section{Results}

\subsection{Random Forest is the strongest standalone model}

\Cref{tab:model-comparison} summarizes every final approach.  The Random Forest--XGBoost blend achieves the best local public mean, 2.191~km.  Random Forest follows at 2.206~km and is the best standalone model; XGBoost is third at 2.276~km.  The same three systems lead on the private and all-trip means.

\begin{table}[H]
  \centering
  \small
  \setlength{\tabcolsep}{5pt}
  \begin{tabular}{@{}rlrrrrr@{}}
\toprule
Rank & Model & Public & Private & All & Validation & Fit (s) \\
\midrule
1 & Ensemble$^{\dagger}$ & 2.191 & 2.500 & 2.346 & 2.113 & -- \\
2 & Random Forest & 2.206 & 2.514 & 2.360 & 2.119 & 25,546.1 \\
3 & XGBoost & 2.276 & 2.527 & 2.402 & 2.215 & 130.4 \\
4 & Decision Tree & 2.375 & 2.732 & 2.553 & 2.328 & 180.3 \\
5 & CatBoost & 2.390 & 2.565 & 2.478 & 2.297 & 473.0 \\
6 & Ridge & 2.665 & 3.107 & 2.886 & 2.888 & 28.1 \\
7 & Linear Regression & 2.668 & 3.099 & 2.883 & 2.891 & 53.1 \\
8 & KNN$^{\ddagger}$ & 3.351 & 3.568 & 3.460 & 3.566 & 0.577 \\
\bottomrule
\end{tabular}

  \vspace{0.35em}

  \begin{minipage}{0.96\linewidth}
    \footnotesize
    $^{\dagger}$The ensemble validation value is the mean of five repeated grouped cross-fits; the blend adds no new base-model fit.
    $^{\ddagger}$KNN validation and fit time use sampled promotion cohorts.  Other standalone validation values use the complete shared holdout.  Recorded fit times span CPU and GPU runs.
  \end{minipage}
  \caption[Final model comparison]{Final model comparison.  Distances are mean Haversine errors in kilometres; lower is better.  Public and private each contain 160 trips.}
  \label{tab:model-comparison}
\end{table}

The public ranking in \cref{fig:public-score-ranking} makes the top-level separation clear: the ensemble is first, Random Forest is the strongest standalone model, and XGBoost is the next-best complementary system.  The ensemble's 15.1~m advantage over Random Forest is nevertheless small relative to their roughly 2.2~km mean errors, so the chart should be read as a ranking rather than evidence of a large practical gap.

\begin{figure}[H]
  \centering
  \begin{minipage}{0.86\linewidth}
    \centering
    \includegraphics[width=\linewidth]{public_score_ranking.png}
    \caption[Local public-score ranking]{Local public-score ranking on the 160 trips designated public.  Bars show mean Haversine error in kilometres; lower is better, and gold highlights the promoted ensemble.}
    \label{fig:public-score-ranking}
  \end{minipage}
\end{figure}

The nonlinear tree families clearly outperform the linear and neighbourhood baselines.  Decision Tree improves the public mean by about 293~m relative to Linear Regression, and bagging reduces it by a further 168~m.  Ridge and Linear Regression are almost tied, showing that regularization does not overcome the missing nonlinear interactions.  KNN ranks last: Manhattan distance in a 541-dimensional mixed spatial and one-hot space is not an effective global notion of route similarity at this scale.

The lower ranks depend on the evaluated subset.  CatBoost is better than Decision Tree on validation, private mean, and all-trip mean, but Decision Tree is 15.8~m better on the public subset.  Ridge is 2.9~m better than Linear Regression publicly, while Linear Regression is slightly better on the private and all-trip means.  These small reversals should not be interpreted as stable superiority between close models.

\subsection{Validation and local scoring tell a consistent top-level story}

On the full standalone holdout, Random Forest scores 2.119~km, followed by XGBoost at 2.215~km, CatBoost at 2.297~km, Decision Tree at 2.328~km, Ridge at 2.888~km, and Linear Regression at 2.891~km.  KNN's 3.566~km value comes from its sampled promotion cohort and is not directly comparable.  The grouped cross-fitted ensemble reaches 2.113~km.  Thus validation and local competition scoring agree on the central conclusion: Random Forest is the strongest base model and XGBoost is the most useful complementary model.

Random Forest's advantage comes with substantial CPU cost.  Its accepted validation refit took 25,546~s, compared with 130~s for GPU XGBoost, 473~s for GPU CatBoost, and 180~s for Decision Tree.  XGBoost therefore provides the best observed accuracy--runtime compromise in this workflow.  The approximately $196\times$ Random Forest/XGBoost ratio is operational rather than algorithmic, because the runs use different hardware.

\subsection{The ensemble gain is consistent but small}

The frozen blend assigns 81.399\% weight to Random Forest and 18.601\% to XGBoost.  It improves on Random Forest in all 25 held-out group folds and in all five repeat means (\cref{fig:ensemble-repeat-gain}).  The mean cross-fitted gain is 6.531~m, with a trip-group robust 95\% interval of [4.555, 8.506]~m.  These results pass the pre-defined promotion gate and support using the blend instead of Random Forest alone.

\begin{figure}[H]
  \centering
  \begin{minipage}{0.84\linewidth}
    \centering
    \includegraphics[width=\linewidth]{ensemble_repeat_gain.png}
    \caption[Mean ensemble gain over Random Forest by cross-fit repeat]{Mean ensemble gain over Random Forest in each repeated five-fold trip-group cross-fit.  Positive values favour the blend; the shaded band is the 95\% trip-group robust interval for the overall mean gain.}
    \label{fig:ensemble-repeat-gain}
  \end{minipage}
\end{figure}

On the local public subset, the ensemble improves the mean by 15.085~m (0.684\%) and ranks first at 2.191~km.  Across all 320 trips it improves the mean by 14.615~m.  The effect is not uniform: Random Forest has the slightly lower all-trip median (1.236 versus 1.257~km), while the ensemble has the lower 90th percentile (6.222 versus 6.243~km).  The blend is therefore an incremental correction that reduces some larger errors rather than a uniformly better prediction on every trip.
